\documentclass[10pt,twocolumn]{article}
% Compile-ready with a stock TeX install:  pdflatex paper.tex
\usepackage[T1]{fontenc}
\usepackage[utf8]{inputenc}
\usepackage{times}
\usepackage{amsmath,amssymb}
\usepackage{booktabs}
\usepackage{graphicx}
\usepackage{url}
\usepackage[margin=0.8in]{geometry}
\usepackage{hyperref}

\title{\textbf{NeuroGraphRAG: Community-Aware Cross-Lingual GraphRAG\\ for Multilingual Neuroscience Retrieval}}
\author{NeuroGraphRAG contributors}
\date{2026}

\begin{document}
\maketitle

\begin{abstract}
Retrieval-augmented generation over scientific literature is dominated by monolingual, flat dense
retrieval, which underperforms on jargon-dense, multilingual, and \emph{global} information needs. We
present \textbf{NeuroGraphRAG}, a graph-augmented retrieval system for neuroscience (EEG/ERP and
biomedical literature) that (i) grounds a knowledge graph, a graph retriever, and a text embedder in a
single curated ontology with \emph{multilingual aliases}, yielding cross-lingual concept embeddings
without any trained model; and (ii) introduces \textbf{Community-Aware Cross-Lingual Retrieval Fusion
(C\textsuperscript{2}RF)}, which augments reciprocal-rank fusion of lexical, dense, and graph signals
with a knowledge-graph community prior; and (iii) a \textbf{Conformal GraphRAG} layer that returns
variable-size retrieval sets with distribution-free coverage guarantees, plus community-conditional and
cross-lingual calibration and risk-controlled abstention. On a bundled multilingual benchmark (59
passages, 72 queries, four languages), C\textsuperscript{2}RF attains the best nDCG@10 (0.722), MAP
(0.608), and Recall@5 (0.633), and wins every query type. The conformal layer's coverage guarantee holds
empirically yet exposes a hidden failure -- the worst knowledge-graph community is covered only 3.7\% of
the time under marginal calibration, repaired 6$\times$ by Mondrian calibration -- and shows that
English-calibrated guarantees under-cover Bengali/Tamil. The pipeline runs offline on CPU; optional
\texttt{sentence-transformers}+LoRA/PEFT and FAISS backends are supported.
\end{abstract}

\section{Introduction}
Neuroscience is a demanding setting for scientific retrieval: terminology is dense and
abbreviation-heavy (N400, MMN, SSVEP, ICA), many readers are more comfortable in Indic languages than in
English, and many real questions are \emph{global} (``which ERP components are altered in
schizophrenia?'') rather than a single-passage lookup. GraphRAG~\cite{edge2024graphrag} added
knowledge-graph structure and community summaries for global questions, but existing pipelines are
English-centric and lean heavily on large LLMs. We ask whether lightweight, ontology-grounded structure
can deliver GraphRAG-style cross-lingual and global benefits without a large model in the loop.

Our contributions: (1) \emph{ontology-grounded concept embeddings} that are cross-lingual by
construction; (2) a precise seed-centric \emph{graph-expansion retriever} fused with BM25 and dense
signals; (3) \emph{C\textsuperscript{2}RF}, a multiplicative community prior over rank fusion; and (4) a
reproducible, offline evaluation harness with a full ablation grid.

\section{Method}
\textbf{Ontology and matching.} Typed concepts carry a canonical form, per-language aliases, and seed
relations. A single matcher tags text with concept ids and is reused by the graph, the graph retriever,
and the embedder, keeping them consistent.

\textbf{Cross-lingual concept embeddings.} A passage maps to
$v=[\,w\cdot\mathrm{onehot}(\mathrm{concepts})\,;\,\mathrm{hash}(\text{char n-grams})\,]$, L2-normalized.
The concept block is language-independent, so parallel English/Hindi passages share dimensions and align.

\textbf{Knowledge graph.} Concept nodes are linked by ontology seed relations and corpus co-occurrence;
weighted greedy-modularity communities each receive a templated summary for global context.

\textbf{Retrievers.} BM25 (Indic-aware tokenization); dense cosine search (FAISS or numpy); and a graph
retriever scoring passages by hop-decayed concept coverage (seed 2.0, 1-hop 0.3, 2-hop 0.1).

\textbf{C\textsuperscript{2}RF.} On top of reciprocal-rank fusion
$s(p)=\sum_r w_r/(k+\mathrm{rank}_r(p))$, we apply $s(p)\leftarrow s(p)\,(1+\lambda\,a(p))$, where the
affinity $a(p)$ is the community-weight mass of the query-relevant concepts that $p$ actually contains
(capped at 1). Gating on concept overlap guarantees an on-topic passage is never demoted below a merely
same-community distractor.

\section{Experiments}
The benchmark has 59 concept-dense passages (English with parallel Hindi/Bengali/Tamil) and 72 typed
queries (factoid, cross-lingual, global, multi-hop) with binary relevance and reference answers. Primary
metric is nDCG@10; backend \texttt{concept-hash}; seed 20260713.

\begin{table}[t]
\centering\small
\caption{Main results on the bundled benchmark.}
\begin{tabular}{lcccc}
\toprule
Configuration & MRR & MAP & R@5 & nDCG@10 \\
\midrule
BM25 & 0.932 & 0.581 & 0.606 & 0.687 \\
Dense & 0.934 & 0.557 & 0.573 & 0.665 \\
Graph & 0.582 & 0.530 & 0.617 & 0.618 \\
RRF (BM25+Dense) & \textbf{0.954} & 0.585 & 0.601 & 0.691 \\
RRF (+Graph) & 0.941 & 0.584 & 0.627 & 0.693 \\
\textbf{C\textsuperscript{2}RF (ours)} & 0.941 & \textbf{0.608} & \textbf{0.633} & \textbf{0.722} \\
\bottomrule
\end{tabular}
\end{table}

\begin{table}[t]
\centering\small
\caption{nDCG@10 by query type.}
\begin{tabular}{lcccc}
\toprule
Configuration & fact. & cross-ling. & global & multi-hop \\
\midrule
BM25 & 0.763 & 0.678 & 0.550 & 0.692 \\
Dense & 0.730 & 0.668 & 0.494 & 0.712 \\
RRF (B+D) & 0.758 & 0.691 & 0.549 & 0.685 \\
\textbf{C\textsuperscript{2}RF} & \textbf{0.785} & \textbf{0.721} & \textbf{0.588} & \textbf{0.735} \\
\bottomrule
\end{tabular}
\end{table}

\textbf{Findings.} Fusion beats any single retriever; the graph signal adds recall (Recall@5
$0.601\!\to\!0.627$); and C\textsuperscript{2}RF lifts nDCG@10 $0.693\!\to\!0.722$, winning every query
type and concentrated on cross-lingual and global queries. On cross-lingual queries, concept-grounded
configurations exceed BM25 (0.721/0.668 vs.\ 0.678) while lexical retrieval collapses on global queries,
isolating the value of shared-concept embeddings and graph context.

\section{Conformal GraphRAG: risk control}
We wrap the retriever in a conformal layer returning a variable-size set
$C(q)=\{d: s(q,d)\ge\tau\}$ with $P(C(q)\cap\mathrm{Rel}(q)\neq\emptyset)\ge 1-\alpha$, where $\tau$ is
the $\lfloor\alpha(n+1)\rfloor$ order statistic of calibration anchors $r(q)=\max_{d\in\mathrm{Rel}(q)}
s(q,d)$; the guarantee is distribution-free under query exchangeability. Empirically coverage tracks the
target (0.814 at $1-\alpha=0.80$) with sets shrinking to $\sim$1 of 59 documents. Crucially, marginal
calibration hides conditional failures: the worst knowledge-graph community is covered only 0.037,
raised $6\times$ to 0.222 by community-conditional (Mondrian) calibration. Calibrating on English
under-covers Bengali (0.48) and Tamil (0.63), a cross-lingual distribution-shift effect that pooled
calibration narrows. Learn-then-Test selective generation bounds answered-query error at a target
$\beta$ (e.g.\ $\beta{=}0.10\to0.083$ at 94\% answer rate).

\section{Limitations and Conclusion}
The benchmark is small and partly synthetic, so the contribution is the relative ordering of
configurations, the conformal guarantees, and their mechanisms. NeuroGraphRAG shows ontology-grounded
structure delivers GraphRAG-style cross-lingual and global benefits cheaply and reproducibly, with
C\textsuperscript{2}RF a small, interpretable addition to rank fusion, and a conformal layer that turns
the retriever into a risk-controlled system exposing failures average metrics hide.

\begin{thebibliography}{9}
\bibitem{edge2024graphrag} D.\ Edge et al.\ From Local to Global: A Graph RAG Approach to Query-Focused Summarization. 2024.
\bibitem{rrf} G.\ Cormack, C.\ Clarke, S.\ Buettcher. Reciprocal Rank Fusion. SIGIR 2009.
\bibitem{dpr} V.\ Karpukhin et al.\ Dense Passage Retrieval. EMNLP 2020.
\bibitem{labse} F.\ Feng et al.\ Language-agnostic BERT Sentence Embedding. 2020.
\bibitem{lora} E.\ Hu et al.\ LoRA: Low-Rank Adaptation of Large Language Models. ICLR 2022.
\bibitem{ragas} S.\ Es et al.\ RAGAS. 2023.
\bibitem{angelopoulos2023} A.\ Angelopoulos, S.\ Bates. A Gentle Introduction to Conformal Prediction. 2023.
\bibitem{ltt} A.\ Angelopoulos et al.\ Learn Then Test: Calibrating Predictive Algorithms to Achieve Risk Control. 2022.
\bibitem{mondrian} V.\ Vovk. Conditional Validity of Inductive Conformal Predictors. ACML 2012.
\end{thebibliography}

\end{document}
